% Options for packages loaded elsewhere
\PassOptionsToPackage{unicode}{hyperref}
\PassOptionsToPackage{hyphens}{url}
%
\documentclass[
  11pt,
]{article}
\usepackage{amsmath,amssymb}
\usepackage{iftex}
\ifPDFTeX
  \usepackage[T1]{fontenc}
  \usepackage[utf8]{inputenc}
  \usepackage{textcomp} % provide euro and other symbols
\else % if luatex or xetex
  \usepackage{unicode-math} % this also loads fontspec
  \defaultfontfeatures{Scale=MatchLowercase}
  \defaultfontfeatures[\rmfamily]{Ligatures=TeX,Scale=1}
\fi
\usepackage{lmodern}
\ifPDFTeX\else
  % xetex/luatex font selection
  \setmainfont[]{DejaVu Sans}
  \setsansfont[]{DejaVu Sans}
  \setmonofont[]{DejaVu Sans Mono}
\fi
% Use upquote if available, for straight quotes in verbatim environments
\IfFileExists{upquote.sty}{\usepackage{upquote}}{}
\IfFileExists{microtype.sty}{% use microtype if available
  \usepackage[]{microtype}
  \UseMicrotypeSet[protrusion]{basicmath} % disable protrusion for tt fonts
}{}
\makeatletter
\@ifundefined{KOMAClassName}{% if non-KOMA class
  \IfFileExists{parskip.sty}{%
    \usepackage{parskip}
  }{% else
    \setlength{\parindent}{0pt}
    \setlength{\parskip}{6pt plus 2pt minus 1pt}}
}{% if KOMA class
  \KOMAoptions{parskip=half}}
\makeatother
\usepackage{xcolor}
\usepackage[a4paper,margin=2.15cm]{geometry}
\usepackage{color}
\usepackage{fancyvrb}
\newcommand{\VerbBar}{|}
\newcommand{\VERB}{\Verb[commandchars=\\\{\}]}
\DefineVerbatimEnvironment{Highlighting}{Verbatim}{commandchars=\\\{\}}
% Add ',fontsize=\small' for more characters per line
\newenvironment{Shaded}{}{}
\newcommand{\AlertTok}[1]{\textcolor[rgb]{1.00,0.00,0.00}{\textbf{#1}}}
\newcommand{\AnnotationTok}[1]{\textcolor[rgb]{0.38,0.63,0.69}{\textbf{\textit{#1}}}}
\newcommand{\AttributeTok}[1]{\textcolor[rgb]{0.49,0.56,0.16}{#1}}
\newcommand{\BaseNTok}[1]{\textcolor[rgb]{0.25,0.63,0.44}{#1}}
\newcommand{\BuiltInTok}[1]{\textcolor[rgb]{0.00,0.50,0.00}{#1}}
\newcommand{\CharTok}[1]{\textcolor[rgb]{0.25,0.44,0.63}{#1}}
\newcommand{\CommentTok}[1]{\textcolor[rgb]{0.38,0.63,0.69}{\textit{#1}}}
\newcommand{\CommentVarTok}[1]{\textcolor[rgb]{0.38,0.63,0.69}{\textbf{\textit{#1}}}}
\newcommand{\ConstantTok}[1]{\textcolor[rgb]{0.53,0.00,0.00}{#1}}
\newcommand{\ControlFlowTok}[1]{\textcolor[rgb]{0.00,0.44,0.13}{\textbf{#1}}}
\newcommand{\DataTypeTok}[1]{\textcolor[rgb]{0.56,0.13,0.00}{#1}}
\newcommand{\DecValTok}[1]{\textcolor[rgb]{0.25,0.63,0.44}{#1}}
\newcommand{\DocumentationTok}[1]{\textcolor[rgb]{0.73,0.13,0.13}{\textit{#1}}}
\newcommand{\ErrorTok}[1]{\textcolor[rgb]{1.00,0.00,0.00}{\textbf{#1}}}
\newcommand{\ExtensionTok}[1]{#1}
\newcommand{\FloatTok}[1]{\textcolor[rgb]{0.25,0.63,0.44}{#1}}
\newcommand{\FunctionTok}[1]{\textcolor[rgb]{0.02,0.16,0.49}{#1}}
\newcommand{\ImportTok}[1]{\textcolor[rgb]{0.00,0.50,0.00}{\textbf{#1}}}
\newcommand{\InformationTok}[1]{\textcolor[rgb]{0.38,0.63,0.69}{\textbf{\textit{#1}}}}
\newcommand{\KeywordTok}[1]{\textcolor[rgb]{0.00,0.44,0.13}{\textbf{#1}}}
\newcommand{\NormalTok}[1]{#1}
\newcommand{\OperatorTok}[1]{\textcolor[rgb]{0.40,0.40,0.40}{#1}}
\newcommand{\OtherTok}[1]{\textcolor[rgb]{0.00,0.44,0.13}{#1}}
\newcommand{\PreprocessorTok}[1]{\textcolor[rgb]{0.74,0.48,0.00}{#1}}
\newcommand{\RegionMarkerTok}[1]{#1}
\newcommand{\SpecialCharTok}[1]{\textcolor[rgb]{0.25,0.44,0.63}{#1}}
\newcommand{\SpecialStringTok}[1]{\textcolor[rgb]{0.73,0.40,0.53}{#1}}
\newcommand{\StringTok}[1]{\textcolor[rgb]{0.25,0.44,0.63}{#1}}
\newcommand{\VariableTok}[1]{\textcolor[rgb]{0.10,0.09,0.49}{#1}}
\newcommand{\VerbatimStringTok}[1]{\textcolor[rgb]{0.25,0.44,0.63}{#1}}
\newcommand{\WarningTok}[1]{\textcolor[rgb]{0.38,0.63,0.69}{\textbf{\textit{#1}}}}
\usepackage{longtable,booktabs,array}
\usepackage{calc} % for calculating minipage widths
% Correct order of tables after \paragraph or \subparagraph
\usepackage{etoolbox}
\makeatletter
\patchcmd\longtable{\par}{\if@noskipsec\mbox{}\fi\par}{}{}
\makeatother
% Allow footnotes in longtable head/foot
\IfFileExists{footnotehyper.sty}{\usepackage{footnotehyper}}{\usepackage{footnote}}
\makesavenoteenv{longtable}
\usepackage{graphicx}
\makeatletter
\def\maxwidth{\ifdim\Gin@nat@width>\linewidth\linewidth\else\Gin@nat@width\fi}
\def\maxheight{\ifdim\Gin@nat@height>\textheight\textheight\else\Gin@nat@height\fi}
\makeatother
% Scale images if necessary, so that they will not overflow the page
% margins by default, and it is still possible to overwrite the defaults
% using explicit options in \includegraphics[width, height, ...]{}
\setkeys{Gin}{width=\maxwidth,height=\maxheight,keepaspectratio}
% Set default figure placement to htbp
\makeatletter
\def\fps@figure{htbp}
\makeatother
\setlength{\emergencystretch}{3em} % prevent overfull lines
\providecommand{\tightlist}{%
  \setlength{\itemsep}{0pt}\setlength{\parskip}{0pt}}
\setcounter{secnumdepth}{5}
\ifLuaTeX
\usepackage[bidi=basic]{babel}
\else
\usepackage[bidi=default]{babel}
\fi
\babelprovide[main,import]{russian}
\ifPDFTeX
\else
\babelfont{rm}[]{DejaVu Sans}
\fi
% get rid of language-specific shorthands (see #6817):
\let\LanguageShortHands\languageshorthands
\def\languageshorthands#1{}
% Custom LaTeX styling for the n-step SAC project report.
% This file is included through: pandoc ... -H paper/latex/preamble.tex
\usepackage{xcolor}
\usepackage{graphicx}
\usepackage{float}
\usepackage{booktabs}
\usepackage{longtable}
\usepackage{array}
\usepackage{caption}
\usepackage{enumitem}
\usepackage{titlesec}
\usepackage{fancyhdr}
\usepackage[most]{tcolorbox}

\definecolor{MTCred}{HTML}{E30613}
\definecolor{DeepText}{HTML}{202124}
\definecolor{SoftGray}{HTML}{666A73}
\definecolor{PaleLavender}{HTML}{F4F0FA}
\definecolor{PaleBlue}{HTML}{EEF6FF}
\definecolor{CodeBg}{HTML}{F6F7F9}
\definecolor{RuleGray}{HTML}{D9DDE7}

\AtBeginDocument{\hypersetup{
  colorlinks=true,
  linkcolor=MTCred,
  urlcolor=MTCred,
  citecolor=MTCred,
  pdfborder={0 0 0}
}}

\captionsetup{font=small, labelfont=bf, labelsep=period}
\setlist[itemize]{leftmargin=1.4em, itemsep=0.25em, topsep=0.25em}
\setlist[enumerate]{leftmargin=1.6em, itemsep=0.25em, topsep=0.25em}

\titleformat{\section}
  {\Large\bfseries\sffamily\color{DeepText}}
  {\thesection}{0.6em}{}
\titleformat{\subsection}
  {\large\bfseries\sffamily\color{DeepText}}
  {\thesubsection}{0.6em}{}
\titlespacing*{\section}{0pt}{1.4em}{0.55em}
\titlespacing*{\subsection}{0pt}{1.0em}{0.35em}

\pagestyle{fancy}
\fancyhf{}
\lhead{\small\sffamily n-step SAC final project}
\rhead{\small\sffamily Radmir I.}
\cfoot{\small\sffamily\thepage}
\renewcommand{\headrulewidth}{0.3pt}
\renewcommand{\footrulewidth}{0pt}

\newtcolorbox{keybox}{
  colback=PaleLavender,
  colframe=MTCred,
  boxrule=0.8pt,
  arc=2mm,
  left=2.5mm,
  right=2.5mm,
  top=2mm,
  bottom=2mm,
  fontupper=\small
}

\newtcolorbox{resultbox}{
  colback=PaleBlue,
  colframe=RuleGray,
  boxrule=0.5pt,
  arc=2mm,
  left=2.5mm,
  right=2.5mm,
  top=2mm,
  bottom=2mm,
  fontupper=\small
}

\newtcolorbox{warningbox}{
  colback=CodeBg,
  colframe=SoftGray,
  boxrule=0.5pt,
  arc=2mm,
  left=2.5mm,
  right=2.5mm,
  top=2mm,
  bottom=2mm,
  fontupper=\small
}

\renewcommand{\arraystretch}{1.18}
\setlength{\parskip}{0.45em}
\setlength{\parindent}{0pt}
\setlength{\emergencystretch}{3em}
\ifLuaTeX
  \usepackage{selnolig}  % disable illegal ligatures
\fi
\usepackage{bookmark}
\IfFileExists{xurl.sty}{\usepackage{xurl}}{} % add URL line breaks if available
\urlstyle{same}
\hypersetup{
  pdftitle={Stress-Testing n-step Soft Actor-Critic under Limited Compute},
  pdfauthor={Radmir I.},
  pdflang={ru-RU},
  hidelinks,
  pdfcreator={LaTeX via pandoc}}

\title{Stress-Testing n-step Soft Actor-Critic under Limited Compute}
\usepackage{etoolbox}
\makeatletter
\providecommand{\subtitle}[1]{% add subtitle to \maketitle
  \apptocmd{\@title}{\par {\large #1 \par}}{}{}
}
\makeatother
\subtitle{Краткий проектный отчет по RL final project}
\author{Radmir I.}
\date{2026}

\begin{document}
\maketitle

\renewcommand*\contentsname{Содержание}
{
\setcounter{tocdepth}{3}
\tableofcontents
}
\begin{keybox}
\textbf{Коротко.} Проект проверяет одну центральную идею из n-step SAC: помогает ли замена стандартного 1-step critic target на n-step target в условиях ограниченного compute. Итоговый результат на \texttt{HalfCheetah-v5}: \texttt{n=3} улучшил обучение относительно baseline, а \texttt{n=5} провалился при тех же hyperparameters. Главный вывод: n-step horizon - это параметр стабильности, а не ручка «чем больше, тем лучше».
\end{keybox}

\section{Abstract}\label{abstract}

This project studies a compact n-step modification of Soft Actor-Critic
(SAC) under a limited compute budget. The main question is whether
replacing the standard 1-step critic target with n-step targets improves
sample efficiency in continuous-control tasks, and where the horizon
becomes unstable. I implement a configurable n-step buffer before the
replay buffer and compare SAC with \texttt{n=1}, \texttt{n=3}, and
\texttt{n=5}. The experiments include \texttt{Pendulum-v1} as a sanity
check and \texttt{HalfCheetah-v5} as the main benchmark. On
\texttt{HalfCheetah-v5}, \texttt{n=3} reaches the best final return,
while \texttt{n=5} fails to learn a strong running policy under the same
hyperparameters. The result suggests that n-step credit assignment can
help SAC, but the horizon is a stability parameter rather than a simple
bigger-is-better knob.

\section{Введение}\label{ux432ux432ux435ux434ux435ux43dux438ux435}

Soft Actor-Critic (SAC) - сильный off-policy actor-critic алгоритм для
задач continuous control. В стандартной постановке critic обучается по
one-step temporal-difference target. Такой target устойчивый и простой,
но распространяет информацию о награде назад только на один шаг.

N-step returns предлагают более агрессивную идею: накопить несколько
rewards перед bootstrap. Это может ускорить credit assignment, особенно
если полезное действие влияет на reward не сразу. Но в off-policy
обучении это не бесплатно: более длинный фрагмент из replay buffer мог
быть собран старой политикой, что повышает off-policy bias и target
variance.

Поэтому проект не пытается заявить новый SOTA-алгоритм. Вместо этого он
изолирует один механизм - horizon critic target - и проверяет его на
небольшом воспроизводимом эксперименте.

\section{Research question}\label{research-question}

Главный вопрос проекта:

\begin{quote}
Помогают ли n-step returns алгоритму SAC при ограниченном compute, и
насколько чувствителен результат к выбранному horizon?
\end{quote}

Сравнивались три варианта:

\begin{longtable}[]{@{}ll@{}}
\toprule\noalign{}
Method & Description \\
\midrule\noalign{}
\endhead
\bottomrule\noalign{}
\endlastfoot
SAC n=1 & standard 1-step SAC baseline \\
SAC n=3 & SAC with moderate 3-step target \\
SAC n=5 & SAC with longer and riskier 5-step target \\
\end{longtable}

Цель - не доказать универсальную оптимальность одного horizon, а
показать, проявляется ли ожидаемый trade-off в компактной и
воспроизводимой постановке.

\section{Контекст
литературы}\label{ux43aux43eux43dux442ux435ux43aux441ux442-ux43bux438ux442ux435ux440ux430ux442ux443ux440ux44b}

Проект находится между стандартным SAC и свежими идеями n-step или
sequence-aware SAC.

\begin{itemize}
\tightlist
\item
  \textbf{SAC} - базовый off-policy actor-critic с entropy
  regularization, twin critics, target networks и replay buffer.
\item
  \textbf{SACn} важен как направление, которое фокусируется на сложности
  n-step returns в off-policy SAC и мотивирует stabilization/correction.
\item
  \textbf{T-SAC / Chunking the Critic} важен тем, что использует
  trajectory chunks и n-step targets вместе с sequence-aware critic.
\end{itemize}

Полное воспроизведение SACn или T-SAC не делалось. Вместо этого проект
тестирует общий для этих направлений механизм: что происходит, когда в
SAC меняется n-step horizon critic target.

\section{Метод}\label{ux43cux435ux442ux43eux434}

Реализация сохраняет стандартную структуру SAC:

\begin{itemize}
\tightlist
\item
  squashed Gaussian actor;
\item
  twin Q-critics;
\item
  target critic networks;
\item
  replay buffer;
\item
  entropy regularization and automatic alpha tuning;
\item
  periodic deterministic evaluation.
\end{itemize}

Главное изменение - \texttt{n-step\ buffer}, который стоит перед replay
buffer. Для окна transitions он строит aggregated transition:

\begin{Shaded}
\begin{Highlighting}[]
\NormalTok{reward\_n  = r\_t + gamma r\_\{t+1\} + ... + gamma\^{}(n{-}1) r\_\{t+n{-}1\}}
\NormalTok{discount\_n = gamma\^{}k}
\NormalTok{next\_state = s\_\{t+k\}}
\end{Highlighting}
\end{Shaded}

где \texttt{k\ \textless{}=\ n}, если эпизод завершился раньше. Critic
target после этого имеет вид:

\begin{Shaded}
\begin{Highlighting}[]
\NormalTok{target = reward\_n + discount\_n * target\_value(next\_state)}
\end{Highlighting}
\end{Shaded}

Для \texttt{n=1} это сводится к обычному SAC.

\begin{resultbox}
\textbf{Почему это хороший ablation?} В экспериментах меняется только horizon critic target: \texttt{n=1}, \texttt{n=3}, \texttt{n=5}. Hyperparameters оставлены одинаковыми, чтобы проверить именно чувствительность к horizon, а не эффект ручного tuning.
\end{resultbox}

\section{Experimental setup}\label{experimental-setup}

Использовались две среды:

\begin{longtable}[]{@{}llrr@{}}
\toprule\noalign{}
Environment & Role & Steps & Seed \\
\midrule\noalign{}
\endhead
\bottomrule\noalign{}
\endlastfoot
\texttt{Pendulum-v1} & sanity check & 30,000 & 0 \\
\texttt{HalfCheetah-v5} & main benchmark & 100,000 & 0 \\
\end{longtable}

Обучение выполнялось на Ryzen 5700 laptop через WSL Ubuntu, CPU
training. Evaluation делался deterministic policy rollout'ами и
сохранялся в CSV. В репозитории включены configs, raw logs, plots и
summary table.

\section{Results}\label{results}

\subsection{Pendulum-v1}\label{pendulum-v1}

На \texttt{Pendulum-v1} все три варианта успешно обучились и пришли к
похожему final quality. Этот эксперимент используется как sanity check:
он подтверждает, что реализация работает и n-step target construction не
ломает SAC на простой среде.

\begin{figure}
\centering
\includegraphics[width=0.92\textwidth,height=\textheight]{results/plots/pendulum_learning_curves.png}
\caption{Pendulum-v1 single-seed n-step SAC ablation}
\end{figure}

\subsection{HalfCheetah-v5}\label{halfcheetah-v5}

Главный эксперимент - ablation на \texttt{HalfCheetah-v5}.

\begin{longtable}[]{@{}lrr@{}}
\toprule\noalign{}
Method & Final eval return & Best eval return \\
\midrule\noalign{}
\endhead
\bottomrule\noalign{}
\endlastfoot
SAC n=1 & 4354.49 & 4354.49 \\
SAC n=3 & 4578.91 & 4578.91 \\
SAC n=5 & 764.49 & 833.37 \\
\end{longtable}

\begin{figure}
\centering
\includegraphics[width=0.92\textwidth,height=\textheight]{results/plots/halfcheetah_learning_curves.png}
\caption{HalfCheetah-v5 single-seed n-step SAC ablation}
\end{figure}

\begin{keybox}
\textbf{Главный результат.} \texttt{n=3} дал лучший final return и быстрее вышел к сильному поведению. \texttt{n=1} остался сильным baseline. \texttt{n=5} не научился устойчивому running policy при тех же hyperparameters.
\end{keybox}

\section{Discussion}\label{discussion}

Результаты поддерживают основную гипотезу: moderate n-step credit
assignment может помочь SAC, но более длинные naive targets могут
destabilize learning.

Самое полезное наблюдение - не только то, что \texttt{n=3} оказался
лучше. Более важен failure mode \texttt{n=5}: при одинаковых
hyperparameters более длинный horizon не улучшил обучение, а почти
сломал его. Это показывает, что horizon нужно воспринимать как stability
parameter.

Сравнение намеренно использует одинаковые hyperparameters для всех
horizons. С одной стороны, это делает ablation чище. С другой стороны,
это может быть не полностью справедливо к \texttt{n=5}, которому могли
бы понадобиться отдельные learning rates, target smoothing или entropy
settings. Но именно это и показывает практическую чувствительность
метода.

\section{Limitations}\label{limitations}

Главное ограничение - single seed. Поэтому проект не должен утверждать,
что \texttt{n=3} всегда лучше \texttt{n=1}, или что \texttt{n=5} всегда
проваливается.

Другие ограничения:

\begin{itemize}
\tightlist
\item
  только две среды;
\item
  нет полного SACn correction mechanism;
\item
  нет transformer или sequence-aware critic;
\item
  нет отдельного hyperparameter tuning для каждого horizon;
\item
  ограниченный CPU compute budget.
\end{itemize}

\begin{warningbox}
\textbf{Корректная формулировка вывода:} в этом limited-compute single-seed ablation, 3-step SAC улучшил HalfCheetah learning относительно 1-step baseline, а 5-step SAC провалился при тех же hyperparameters.
\end{warningbox}

\section{Future work}\label{future-work}

Наиболее полезные продолжения:

\begin{enumerate}
\def\labelenumi{\arabic{enumi}.}
\tightlist
\item
  запустить 3-5 seeds;
\item
  добавить Hopper или Walker2d;
\item
  посчитать confidence intervals и failure rates;
\item
  сравнить naive n-step SAC с SACn-style correction;
\item
  отдельно tune hyperparameters для каждого horizon;
\item
  попробовать sequence-aware critic в духе T-SAC.
\end{enumerate}

\section{Conclusion}\label{conclusion}

Проект реализует и тестирует configurable n-step targets в SAC.
\texttt{Pendulum-v1} подтверждает корректность реализации, а
\texttt{HalfCheetah-v5} показывает основной trade-off: \texttt{n=3}
помогает в этой постановке, но \texttt{n=5} проваливается при тех же
hyperparameters.

Итоговая мысль:

\begin{quote}
N-step horizon is a stability parameter, not a bigger-is-better knob.
\end{quote}

\section{References and project
artifacts}\label{references-and-project-artifacts}

\begin{itemize}
\tightlist
\item
  Repository: \texttt{Therad445/nstep-sac-reproducibility}.
\item
  Main implementation: \texttt{src/rl/nstep\_buffer.py},
  \texttt{src/rl/sac.py}, \texttt{src/train.py}.
\item
  Experiment logs: \texttt{results/raw/*.csv}.
\item
  Plots: \texttt{results/plots/*.png}.
\item
  Summary table: \texttt{results/halfcheetah\_summary.csv}.
\item
  Slides: \texttt{slides/nstep\_sac\_template\_style\_deck\_ru.pptx} and
  \texttt{slides/nstep\_sac\_template\_style\_deck.pptx}.
\end{itemize}

\end{document}
